\section{Dogma central de la biología}

El dogma central de la biología molecular establece el flujo de información genética:

\textbf{ADN → ARN → Proteína}

Este principio fue propuesto por Francis Crick en 1958 y explica cómo la información almacenada en el ADN se transcribe a ARN y luego se traduce a proteínas.

\section{Replicación}

La replicación es el proceso de duplicación del ADN.
\begin{itemize}
    \item Su objetivo es la conservación de la información genética
    \item Proceso semiconservativo: Cada nueva molécula contiene una hebra original y una nueva
    \item Participan enzimas como ADN polimerasa, helicasa y ligasa
    \item Ocurre durante la fase S del ciclo celular
\end{itemize}

\section{Transcripción}

La transcripción es el proceso por el cual se genera una copia de ARNm para que luego genere una proteína.

\subsection{Maduración del ARNm}
\begin{itemize}
    \item Se retiran intrones (partes no codificantes)
    \item Se empalman los exones (partes codificantes)
    \item Se añade capa 5' y cola poli-A
\end{itemize}

\section{Traducción}

La traducción es el proceso de traducir una secuencia de ARNm a una secuencia de aminoácidos.
\begin{itemize}
    \item Este proceso es realizado por los ribosomas
    \item Es guiado por codones (ARNm) y anticodones (ARNt)
    \item Ocurre en el citoplasma
\end{itemize}

\subsection{Etapas de la traducción}
\begin{enumerate}
    \item Iniciación: El ribosoma se ensambla en el ARNm
    \item Elongación: Adición secuencial de aminoácidos
    \item Terminación: Liberación de la proteína cuando se encuentra un codón de stop
\end{enumerate}